\chapter{Bootstrapping and Initialization}
\label{chap:init}

Every operating system has a moment of genesis: the instant at which a passive
collection of bytes on disk becomes a living, running system. In AULinux that
moment is owned by \file{au-init}, a small, statically linked C program that
becomes process ID~1 (PID~1) and never exits for as long as the machine is
powered on. This chapter dissects \file{init/src/init.c} and explains not only
\emph{what} it does, but \emph{why} a PID~1 process must do those things, and
what the current implementation does well or leaves for later.

\section{The Boot Chain: GRUB \texorpdfstring{$\rightarrow$}{->} Kernel \texorpdfstring{$\rightarrow$}{->} au-init}

Before a single line of \file{init.c} executes, two earlier actors must hand
control along a precise chain. AULinux uses the GRUB2 bootloader, whose
configuration is deliberately minimal:

\begin{lstlisting}[language=bash, caption={The complete GRUB configuration (bootloader/grub.cfg)}]
set timeout=5
set default=0

menuentry "AULinux" {
    linux /boot/vmlinuz root=/dev/sda1 rw quiet init=/sbin/init
    initrd /boot/initramfs.img
}
\end{lstlisting}

The boot proceeds in three stages:

\begin{enumerate}
  \item \textbf{GRUB} loads the compressed Linux kernel image
        (\file{/boot/vmlinuz}) and an initial RAM filesystem
        (\file{/boot/initramfs.img}) into memory, then transfers control to the
        kernel. The \code{root=/dev/sda1} parameter tells the kernel which block
        device holds the real root filesystem; \code{rw} requests it be mounted
        read--write; \code{quiet} suppresses verbose kernel logging.
  \item \textbf{The Linux kernel} initializes hardware, mounts the root
        filesystem, and then looks for the very first userspace program to run.
        The crucial parameter here is \code{init=/sbin/init}, which explicitly
        names \file{au-init} (installed as \file{/sbin/init}) as that first
        program. The kernel \code{exec()}s it with PID~1.
  \item \textbf{au-init} takes over as PID~1 and performs every step described
        in the remainder of this chapter.
\end{enumerate}

\begin{infobox}[Why PID 1 is special]
The kernel assigns PID~1 to the first userspace process and grants it two
permanent, non-negotiable responsibilities. First, PID~1 is the
\emph{reaper of last resort}: when any process's parent dies, the orphaned
children are re-parented to PID~1, which must call \code{wait()} on them to
release their kernel resources. Second, PID~1 is \emph{unkillable by default} and
\emph{may never exit} --- if it does, the kernel panics with the famous
``Attempted to kill init!'' message. Every design decision in \file{init.c}
flows from these two facts.
\end{infobox}

\subsection{Why static linking is mandatory}

The header comment of \file{init.c} states the constraint plainly:
``Built with \code{-static} to run before shared libraries are available.''
When the kernel executes PID~1, the dynamic linker (\file{ld.so}), the C library
shared objects, and the filesystems that hold them may not yet be available or
mounted. A dynamically linked \file{au-init} could fail to start simply because
\file{libc.so} could not be located --- a fatal bootstrapping paradox, since the
very program responsible for mounting filesystems would itself depend on a
mounted filesystem. By compiling \file{au-init} as a single self-contained
static binary, every symbol it needs is baked directly into the executable, and
it can run with nothing but the kernel beneath it. We return to the exact
build flags in Section~\ref{sec:init-build}.

\section{Anatomy of au-init: The \texorpdfstring{\code{main()}}{main()} Flow}

The \code{main()} function (\file{init.c}, lines 681--792) is the spine of the
program. Stripped to its phases, it reads as a clean, linear bootstrap sequence
followed by an infinite supervision loop:

\begin{lstlisting}[language=C, caption={The phased structure of main() (abridged)}]
int main(int argc, char *argv[])
{
    /* 1. Handle --version / --help and exit early */
    /* 2. Discover our own PID */
    my_pid = getpid();

    /* 3. Export runtime environment (GOGC, GOMEMLIMIT, PATH) */
    setenv("GOGC", "50", 1);
    setenv("GOMEMLIMIT", "512MiB", 1);
    setenv("PATH", "/bin:/sbin:/usr/bin:/usr/sbin", 1);

    /* 4. Print the boot banner */
    /* 5. Install signal handlers */
    setup_signals();
    /* 6. (PID 1 only) wire up the console */
    if (my_pid == 1) setup_console();
    /* 7. (PID 1 only) mount filesystems, then set hostname */
    if (my_pid == 1) { mount_filesystems(); setup_hostname(); }
    /* 8. (PID 1 only) load + launch configured services */
    if (my_pid == 1) load_services();
    /* 9. Launch the interactive login shell */
    start_shell();
    /* 10. Enter the supervision loop until a shutdown signal */
    while (running) { ... }
    /* 11. Reboot or power off */
}
\end{lstlisting}

A noteworthy design choice is that nearly every privileged action is guarded by
\code{if (my\_pid == 1)}. When the program detects it is \emph{not} PID~1, it
prints a warning and announces ``Running in test mode'' (lines 713--718). This
lets a developer execute the binary from an ordinary shell during development
without it attempting to mount system filesystems or set the hostname --- it
will only try to start a shell. This dual-mode behaviour makes the init binary
testable without a virtual machine.

\subsection{Environment preparation}

Immediately before doing any real work, \file{au-init} exports three environment
variables (lines 700--702) that every descendant process will inherit:

\begin{itemize}
  \item \code{GOGC=50} --- lowers the Go garbage collector's growth threshold so
        the Go-based userland (the shell and utilities) collects more
        aggressively, trading some CPU for a smaller memory footprint.
  \item \code{GOMEMLIMIT=512MiB} --- sets a soft ceiling on Go heap memory.
  \item \code{PATH=/bin:/sbin:/usr/bin:/usr/sbin} --- the canonical system search
        path, so children can resolve command names.
\end{itemize}

\begin{notebox}[A C init tuning a Go userland]
The presence of \code{GOGC} and \code{GOMEMLIMIT} is a direct fingerprint of
AULinux's polyglot design: a C PID~1 deliberately tunes the Go runtime of the
shell and core utilities it is about to spawn. This is unusual for a traditional
init and reflects the project's mixed-language architecture described in
Chapter~1.
\end{notebox}

\section{Mounting the Pseudo-Filesystems}

The kernel hands \file{au-init} a root filesystem but almost nothing else. The
\code{mount\_filesystems()} routine (lines 109--188) populates the standard set
of \emph{pseudo}- and \emph{virtual} filesystems that the rest of userspace
assumes exist. Each mount follows the same defensive pattern: attempt the mount,
and if it fails for any reason other than \code{EBUSY} (already mounted), log a
warning but \emph{keep going}. Treating \code{EBUSY} as success makes the routine
idempotent --- safe to run even if an initramfs already mounted some of these.

\begin{lstlisting}[language=C, caption={Representative mount calls from mount\_filesystems()}]
/* Mount /proc */
if (mount("proc", "/proc", "proc",
          MS_NOSUID | MS_NODEV | MS_NOEXEC, NULL) < 0) {
    if (errno != EBUSY)
        log_msg("WARN", "Failed to mount /proc: %s", strerror(errno));
} else {
    log_msg("INFO", "Mounted /proc");
}
...
/* Mount /dev (devtmpfs) */
if (mount("devtmpfs", "/dev", "devtmpfs", MS_NOSUID, "mode=0755") < 0) {
    if (errno != EBUSY)
        log_msg("WARN", "Failed to mount /dev: %s", strerror(errno));
}
\end{lstlisting}

The complete set of filesystems mounted, in order, is given in
Table~\ref{tab:init-mounts}. The ordering matters: \file{/dev} (a devtmpfs) must
exist before \file{/dev/pts} can be mounted on top of it, which is why
\code{mkdir("/dev/pts", 0755)} (line 141) precedes the devpts mount.

\begin{table}[h]
\centering
\caption{Filesystems mounted by \code{mount\_filesystems()}, in order.}
\label{tab:init-mounts}
\small
\begin{tabular}{@{}llll@{}}
\toprule
\textbf{Source} & \textbf{Target} & \textbf{Type} & \textbf{Purpose} \\
\midrule
proc     & /proc    & proc     & Process and kernel state interface \\
sysfs    & /sys     & sysfs    & Device and kernel object hierarchy \\
devtmpfs & /dev     & devtmpfs & Kernel-managed device nodes \\
devpts   & /dev/pts & devpts   & Pseudo-terminal slave devices \\
tmpfs    & /run     & tmpfs    & Volatile runtime state (PID files etc.) \\
tmpfs    & /tmp     & tmpfs    & Temporary scratch space \\
\bottomrule
\end{tabular}
\end{table}

The mount flags are chosen with security in mind. \code{MS\_NOSUID} disables
set-user-ID bits, \code{MS\_NODEV} forbids device-node interpretation, and
\code{MS\_NOEXEC} forbids executing binaries from the mount. The \file{/proc} and
\file{/sys} mounts carry all three; \file{/dev/pts} is mounted
\code{nosuid,noexec} with \code{gid=5,mode=620} (the conventional \code{tty}
group and permissions for pseudo-terminals); \file{/run} and \file{/tmp} are
\code{nosuid,nodev} tmpfs mounts, with \file{/tmp} given the sticky-bit mode
\code{1777} so any user may create files but only delete their own.

\subsection{Bringing up the loopback interface}

The mount routine also performs one piece of network setup that is easy to
overlook: it brings the loopback interface (\code{lo}) up. Many local services
and IPC mechanisms assume \code{127.0.0.1} is reachable, so init does this
unconditionally using raw \code{ioctl} socket calls rather than shelling out to
an external tool (lines 168--185):

\begin{lstlisting}[language=C, caption={Activating the loopback interface via ioctl}]
int sock = socket(AF_INET, SOCK_DGRAM, 0);
if (sock >= 0) {
    struct ifreq ifr;
    memset(&ifr, 0, sizeof(ifr));
    strncpy(ifr.ifr_name, "lo", IFNAMSIZ);
    if (ioctl(sock, SIOCGIFFLAGS, &ifr) >= 0) {
        ifr.ifr_flags |= (IFF_UP | IFF_RUNNING);
        if (ioctl(sock, SIOCSIFFLAGS, &ifr) < 0)
            log_msg("WARN", "Failed to bring up lo: %s", strerror(errno));
        else
            log_msg("INFO", "Interface lo is up");
    }
    close(sock);
}
\end{lstlisting}

It reads the current interface flags with \code{SIOCGIFFLAGS}, ORs in
\code{IFF\_UP | IFF\_RUNNING}, and writes them back with \code{SIOCSIFFLAGS} ---
exactly what \code{ifconfig lo up} would do, but with no external dependency,
honouring the project's ``no external libraries'' rule for PID~1.

\subsection{Setting the hostname}

Immediately after \code{mount\_filesystems()} returns, \code{main()} calls
\code{setup\_hostname()} (lines 193--211). This dedicated function reads
\file{/etc/hostname}, strips any trailing carriage-return or newline with
\code{strcspn}, and applies the result via the \code{sethostname()} syscall. If
the file is absent or unreadable, the function falls back silently to the
compiled-in default string \code{"aulinux"}, ensuring the system always has a
valid hostname regardless of rootfs state. A \code{log\_msg("INFO",...)} confirms
success; a \code{log\_msg("WARN",...)} records any \code{sethostname()} failure:

\begin{lstlisting}[language=C, caption={setup\_hostname(): reading /etc/hostname with fallback}]
static void setup_hostname(void)
{
    FILE *fp = fopen("/etc/hostname", "r");
    char hostname[128] = "aulinux";       /* safe default */

    if (fp) {
        if (fgets(hostname, sizeof(hostname), fp))
            hostname[strcspn(hostname, "\r\n")] = '\0';
        fclose(fp);
    }

    if (sethostname(hostname, strlen(hostname)) < 0)
        log_msg("WARN", "Failed to set hostname: %s", strerror(errno));
    else
        log_msg("INFO", "Hostname set to '%s'", hostname);
}
\end{lstlisting}

The call sequence in \code{main()} is therefore: signals $\rightarrow$ console
$\rightarrow$ mount filesystems $\rightarrow$ set hostname $\rightarrow$ load
services $\rightarrow$ start shell. Hostname configuration is deliberately placed
after mounting because \file{/etc/hostname} lives on the root filesystem, which
must be fully accessible before the file can be opened.

\subsection{Console setup}

\code{setup\_console()} (lines 216--247) opens \file{/dev/console} and on failure
falls back to \file{/dev/tty1}, then uses \code{dup2()} to wire that descriptor
onto \code{stdin}, \code{stdout}, and \code{stderr} so that init's own log output
and the shell it later spawns are visible on the physical console. Note that
console setup runs \emph{before} the filesystem mounts --- it relies only on the
device nodes provided by the kernel itself during early boot.

\section{Service Configuration and Parsing}

AULinux services are declared in a flat configuration file located at
\file{/etc/aulinux/services.conf} (the \code{SERVICE\_CONFIG} constant, line~36).
The format is intentionally trivial --- one service per line, in the form:

\begin{lstlisting}[language=bash, caption={The services.conf line format}]
# Comments and blank lines are ignored
name=command and its arguments
\end{lstlisting}

The \code{load\_services()} function (lines 467--538) is now divided into two
distinct responsibilities: first it populates the \code{services[]} table, then
it delegates process creation to \code{start\_service()}. The parser reads the
file line by line:

\begin{lstlisting}[language=C, caption={Parsing a service definition line}]
while (fgets(line, sizeof(line), fp)) {
    line[strcspn(line, "\n")] = 0;          /* strip newline      */
    if (line[0] == '#' || line[0] == 0)      /* skip comment/blank */
        continue;

    name = strtok(line, "=");                /* text before '='    */
    cmd  = strtok(NULL, "");                 /* everything after   */
    if (!name || !cmd) continue;

    while (*cmd == ' ') cmd++;               /* trim leading space */
    if (*cmd == '\0') continue;
    ...
}
\end{lstlisting}

For each valid line it populates a \code{struct service} in the global
\code{services[]} array (capacity \code{MAX\_SERVICES = 32}). It stores the full
command string in \code{command\_line}, and separately extracts just the
executable name (the first whitespace-delimited token) into \code{command} for
logging. All bookkeeping fields are explicitly zeroed: \code{pid = 0},
\code{state = SERVICE\_STOPPED}, \code{restart\_count = 0}, \code{last\_start = 0},
\code{next\_restart = 0}. After the whole file is parsed, a second pass calls
\code{start\_service(i)} for every entry (lines 534--537).

The \code{struct service} record (lines 50--59) carries everything the
supervisor needs to track a service's lifecycle:

\begin{lstlisting}[language=C, caption={The service control block}]
struct service {
    char name[64];
    char command[256];
    char command_line[512];   /* full command string incl. arguments */
    pid_t pid;
    enum service_state state; /* STOPPED, STARTING, RUNNING,
                                 STOPPING, FAILED */
    int restart_count;
    time_t last_start;
    time_t next_restart;      /* wall-clock time for next restart */
};
\end{lstlisting}

The \code{command\_line} field (512 bytes) stores the raw command string exactly
as it appeared on the right-hand side of the \code{=} in \file{services.conf},
including all arguments. This is distinct from \code{command} (256 bytes), which
holds only the executable path used for log messages. The \code{next\_restart}
field is a \code{time\_t} wall-clock deadline set by \code{reap\_children()} when
a service exits; the main supervision loop compares it against \code{time(NULL)}
on every tick to decide when to respawn.

\begin{notebox}[The shipped services.conf is empty]
In the current AULinux tree the configuration file at
\file{rootfs/etc/aulinux/services.conf} exists but is completely empty (zero
bytes). Consequently, on a stock boot \code{load\_services()} parses no entries
and the only long-running process init launches is the interactive shell
(Section~\ref{sec:init-shell}). The service-supervision machinery is fully
implemented and ready, but no background services are declared yet. If the file
were missing entirely, init logs a warning (``Could not open service config'')
and continues gracefully.
\end{notebox}

\subsection{Starting a service: \texttt{start\_service()}}

The separation of \emph{configuration loading} from \emph{process spawning} is
a key design improvement in the current codebase. \code{load\_services()} no
longer contains any fork/exec logic; all of that lives in \code{start\_service()}
(lines 425--462), which is also called by the supervision loop when a
previously-stopped service becomes due for restart.

\code{start\_service(int idx)} takes the index of an already-populated entry in
\code{services[]} and turns it into a running process:

\begin{enumerate}
  \item It copies \code{services[idx].command\_line} into a local buffer
        (because \code{strtok} is destructive) and tokenises it on spaces into
        an \code{argv[]} array of up to 63 elements, with a \code{NULL}
        sentinel.
  \item It sets \code{state = SERVICE\_STARTING} and records
        \code{last\_start = time(NULL)}.
  \item It calls \code{spawn\_process(argv[0], argv, services[idx].name)}.
  \item On success it records the returned PID, sets \code{state =
        SERVICE\_RUNNING}, and clears \code{next\_restart = 0}. On failure it
        sets \code{state = SERVICE\_FAILED}.
\end{enumerate}

\subsection{Spawning a process: \texttt{spawn\_process()}}

\code{spawn\_process()} (lines 380--420) is the lowest-level primitive. Its
signature now carries a \code{service\_name} parameter:

\begin{lstlisting}[language=C, caption={spawn\_process() signature}]
static pid_t spawn_process(const char *path,
                           char *const argv[],
                           const char *service_name);
\end{lstlisting}

\begin{lstlisting}[language=C, caption={fork/exec with session and log redirection}]
pid = fork();
if (pid == 0) {                       /* child */
    setsid();                          /* new session, detach from tty */

    if (service_name && strcmp(service_name, "shell") != 0) {
        mkdir("/var", 0755);
        mkdir("/var/log", 0755);
        char log_path[256];
        snprintf(log_path, sizeof(log_path),
                 "/var/log/%s.log", service_name);
        int log_fd = open(log_path,
                          O_WRONLY | O_CREAT | O_APPEND, 0644);
        if (log_fd >= 0) {
            dup2(log_fd, STDOUT_FILENO);
            dup2(log_fd, STDERR_FILENO);
            close(log_fd);
        }
    }
    execv(path, argv);                 /* never returns on success */
    fprintf(stderr, "au-init: execv(%s) failed: %s\n",
            path, strerror(errno));
    _exit(127);
}
return pid;                            /* parent returns child pid */
\end{lstlisting}

Three details are worth highlighting. First, the child calls \code{setsid()} to
become a session leader, detaching it from init's controlling terminal so that a
crash or signal on one service cannot ripple through the session. Second, every
service \emph{except} the shell has its \code{stdout} and \code{stderr}
redirected to \file{/var/log/<name>.log}, giving each background service its own
append-only log; the shell is exempted because it needs the live console. The
\code{mkdir("/var", 0755)} and \code{mkdir("/var/log", 0755)} calls before
opening the log file ensure the directory tree exists even on a minimal rootfs
--- they are no-ops if the directories are already present. Third, if
\code{execv} fails the child exits with status \code{127} (the conventional
``command not found'' code) via \code{\_exit()} --- crucially the underscore
variant, to avoid flushing the parent's \code{stdio} buffers in the forked
child.

\section{Process Supervision: Reaping, Restarts, and Signals}

Once the bootstrap phases complete, \file{au-init} settles into its real job for
the rest of the system's uptime: supervising children. This rests on three
pillars --- signal handling, zombie reaping, and a restart policy.

\subsection{Signal handling}

\code{setup\_signals()} (lines 279--297) installs a single \code{signal\_handler}
for the meaningful signals and explicitly ignores two others:

\begin{lstlisting}[language=C, caption={Signal disposition installed by setup\_signals()}]
struct sigaction sa;
memset(&sa, 0, sizeof(sa));
sa.sa_handler = signal_handler;
sigemptyset(&sa.sa_mask);
sa.sa_flags = SA_RESTART;

sigaction(SIGCHLD, &sa, NULL);
sigaction(SIGTERM, &sa, NULL);
sigaction(SIGINT,  &sa, NULL);
sigaction(SIGUSR1, &sa, NULL);  /* Reboot   */
sigaction(SIGUSR2, &sa, NULL);  /* Poweroff */

signal(SIGHUP,  SIG_IGN);
signal(SIGPIPE, SIG_IGN);
\end{lstlisting}

The handler itself (lines 252--274) is a model of async-signal-safe design: it
does no real work, only sets \code{volatile sig\_atomic\_t} flags that the main
loop later inspects.

\begin{itemize}
  \item \code{SIGCHLD} $\rightarrow$ sets \code{child\_exited = 1}, signalling
        that there are zombies to reap.
  \item \code{SIGTERM} / \code{SIGINT} $\rightarrow$ clears \code{running} and
        requests a power-off.
  \item \code{SIGUSR1} $\rightarrow$ requests a reboot.
  \item \code{SIGUSR2} $\rightarrow$ requests a power-off.
\end{itemize}

\begin{infobox}[Signals as the shutdown API]
AULinux has no separate \code{reboot} or \code{poweroff} daemon. Instead,
\emph{sending a signal to PID~1 is the shutdown interface}: \code{kill -USR1 1}
reboots the machine and \code{kill -USR2 1} powers it off. The \code{--help} text
(lines 667--676) documents exactly this contract. \code{SIGHUP} and
\code{SIGPIPE} are ignored so that a hangup on the console or a broken pipe to a
dying child can never accidentally take down init.
\end{infobox}

\subsection{Reaping zombies}

When a child terminates, the kernel keeps a tiny ``zombie'' entry until the
parent retrieves its exit status. As PID~1, \file{au-init} inherits every
orphaned process in the system, so it must reap diligently or leak the process
table. \code{reap\_children()} (lines 302--375) drains all available zombies in a
non-blocking loop:

\begin{lstlisting}[language=C, caption={The non-blocking reaping loop}]
while ((pid = waitpid(-1, &status, WNOHANG)) > 0) {
    for (i = 0; i < service_count; i++) {
        if (services[i].pid == pid) {
            /* classify exit and schedule restart ... */
            break;
        }
    }
}
child_exited = 0;
\end{lstlisting}

The \code{WNOHANG} flag means \code{waitpid} returns immediately when no more
children have exited, so the \code{while} loop reaps every pending zombie in one
pass and then stops --- handling the case where several children die nearly
simultaneously but only one \code{SIGCHLD} is delivered. Orphaned processes that
init did not itself spawn fall through the inner loop unmatched: they are still
reaped (their kernel slot freed), simply without a restart decision. This is
precisely the ``reaper of last resort'' duty of PID~1.

\subsection{The restart policy}

For tracked services, \code{reap\_children()} first classifies how the process
exited, then decides what to do next. Exit classification uses the standard
\code{WIFEXITED}/\code{WIFSIGNALED} macros. A non-zero exit code sets the
\code{crash} flag; so does termination by an uncaught signal, which additionally
logs the signal name using \code{strsignal(signal\_num)} for human-readable
diagnostics.

A \emph{stability window} check runs before any restart decision: if the service
ran for at least 30 seconds before dying, \code{restart\_count} is reset to zero
(lines 341--343). This prevents a service that normally runs for a long time but
occasionally crashes from accumulating backoff penalties indefinitely --- a
single crash after 30+ seconds of uptime is treated as a fresh start.

The policy then distinguishes three cases:

\begin{enumerate}
  \item \textbf{The interactive shell} (\code{strcmp(name, "shell") == 0}) is
        always restarted immediately regardless of how it exited.
        \code{restart\_count} is reset to zero, \code{next\_restart} is set to
        \code{now} (not \code{now + N}), and \code{state} is set to
        \code{SERVICE\_STOPPED} so the main loop restarts it on the very next
        tick. A dead console should reappear at once.

  \item \textbf{A clean exit} (exit code 0) of a non-shell service resets
        \code{restart\_count} to zero and sets \code{next\_restart = now + 1},
        scheduling a restart one second later. This one-second gap prevents a
        tight respawn loop when a service exits successfully but immediately.
        State is set to \code{SERVICE\_STOPPED}.

  \item \textbf{A crash} --- non-zero exit code, or termination by an uncaught
        signal --- increments \code{restart\_count} and applies
        \emph{exponential backoff} via \code{next\_restart}. State is set to
        \code{SERVICE\_FAILED}.
\end{enumerate}

The backoff calculation (lines 351--360) is the heart of the crash policy:

\begin{lstlisting}[language=C, caption={Exponential backoff on repeated crashes}]
if (crash) {
    services[i].restart_count++;
    int delay = 1 << (services[i].restart_count - 1);  /* 1,2,4,8,... */
    if (delay > 60) delay = 60;                        /* cap at 60s  */
    services[i].next_restart = now + delay;
    services[i].state = SERVICE_FAILED;
    log_msg("WARN", "Service '%s' crashed! Exponential backoff: "
            "restarting in %d seconds (retry count: %d)",
            services[i].name, delay, services[i].restart_count);
}
\end{lstlisting}

The delay doubles with each successive crash (1, 2, 4, 8, 16, 32 seconds) and is
capped at 60 seconds, so a service stuck in a crash loop backs off gracefully
instead of pegging the CPU. The \code{next\_restart} field stores the absolute
wall-clock time at which the main loop should attempt the next \code{start\_service()}
call --- not a countdown, but a deadline that the supervision loop compares
against \code{time(NULL)} on every iteration.

\subsection{The main supervision loop}

All of these mechanisms come together in the main loop (lines 756--781):

\begin{lstlisting}[language=C, caption={The init main loop}]
while (running) {
    struct timeval tv = { .tv_sec = 1, .tv_usec = 0 };
    select(0, NULL, NULL, NULL, &tv);   /* sleep up to 1s, signal-aware */

    if (child_exited)
        reap_children();

    time_t now = time(NULL);
    for (int i = 0; i < service_count; i++) {
        if (services[i].pid == 0 &&
            services[i].next_restart > 0 &&
            now >= services[i].next_restart) {
            if (strcmp(services[i].name, "shell") == 0)
                start_shell();
            else
                start_service(i);
        }
    }
}
\end{lstlisting}

The loop uses \code{select()} with all file-descriptor sets \code{NULL} purely as
a one-second sleep that is interruptible by signals --- when \code{SIGCHLD}
arrives, \code{select} returns early, and the loop reaps immediately rather than
waiting out the full second. Each iteration then performs a clock-driven sweep:
any service whose \code{pid} is 0 (not running) and whose \code{next\_restart}
deadline has passed is restarted. This polling design keeps init's logic simple
and entirely free of race conditions between signal delivery and bookkeeping ---
the handler only sets flags, and the loop owns all state changes.

\subsection{The interactive shell}
\label{sec:init-shell}

\code{start\_shell()} (lines 543--597) launches AULinux's login shell. It first
checks whether the preferred shell \file{/bin/aush} is executable; if not, it
falls back to \file{/bin/sh}. It spawns the shell with the arguments
\code{aush -l} (a login shell) and then registers it in the \code{services[]}
array under the special name \code{"shell"} so the supervision loop will restart
it like any other service:

\begin{lstlisting}[language=C, caption={Selecting and launching the login shell}]
char *shell = DEFAULT_SHELL;            /* "/bin/aush" */
char *argv[] = { "aush", "-l", NULL };

if (access(shell, X_OK) != 0) {         /* aush missing? */
    shell = FALLBACK_SHELL;             /* "/bin/sh"    */
    argv[0] = "sh";
    if (access(shell, X_OK) != 0) {
        log_msg("ERROR", "No shell available");
        return -1;
    }
}
pid = spawn_process(shell, argv, "shell");
\end{lstlisting}

If even this fails, \code{main()} refuses to give up: because a PID~1 that
returns would panic the kernel, it spawns a bare \file{/bin/sh} as an emergency
shell instead (lines 744--748). This guarantees the operator always has \emph{a}
prompt to recover from.

\section{Shutdown and Reboot}

When a shutdown or reboot signal sets \code{running} to 0, the main loop exits
and \code{main()} dispatches to \code{shutdown\_system()} (lines 602--653) with a
flag selecting reboot versus power-off. The teardown sequence is textbook and
deliberately ordered:

\begin{enumerate}
  \item Send \code{SIGTERM} to every tracked service, giving each a chance to
        exit cleanly.
  \item \code{sleep(2)} to allow graceful shutdown.
  \item Send \code{SIGKILL} to any service still alive.
  \item Call \code{reap\_children()} to collect the stragglers.
  \item \code{sync()} to flush the page cache to disk.
  \item \code{umount2()} each pseudo-filesystem in reverse mount order, using
        \code{MNT\_DETACH} for a lazy unmount that succeeds even if the mount is
        momentarily busy.
  \item \code{sync()} once more, then call \code{reboot(RB\_AUTOBOOT)} or
        \code{reboot(RB\_POWER\_OFF)}.
\end{enumerate}

\begin{lstlisting}[language=C, caption={Graceful-then-forceful service teardown}]
/* Send SIGTERM to all services */
for (i = 0; i < service_count; i++)
    if (services[i].pid > 0)
        kill(services[i].pid, SIGTERM);

sleep(2);                              /* grace period */

/* Send SIGKILL to any survivors */
for (i = 0; i < service_count; i++)
    if (services[i].pid > 0)
        kill(services[i].pid, SIGKILL);
\end{lstlisting}

The double \code{sync()} bracketing the unmounts is a classic safety measure:
the first flushes dirty data before filesystems are torn down, the second
catches anything written during teardown, minimising the chance of corruption on
the next boot.

\section{Building and Linking au-init}
\label{sec:init-build}

The build is governed by a compact Makefile in \file{init/}:

\begin{lstlisting}[language=bash, caption={The init Makefile (essential rules)}]
CC = gcc
CFLAGS = -Wall -Wextra -O2 -static

SRC_DIR = src
BUILD_DIR = ../build/init
TARGET = $(BUILD_DIR)/au-init

$(TARGET): $(SRC_DIR)/init.c
	$(CC) $(CFLAGS) $< -o $@
\end{lstlisting}

The flags encode the project's hard requirements:

\begin{itemize}
  \item \code{-static} --- the non-negotiable flag discussed earlier. It links
        the C library and every other dependency directly into \file{au-init},
        producing a self-contained binary that runs before any shared library is
        available.
  \item \code{-Wall -Wextra} --- enables a broad set of compiler warnings, a
        sensible discipline for code that runs as PID~1 where a latent bug is
        catastrophic.
  \item \code{-O2} --- standard optimisation.
\end{itemize}

The binary is emitted as \file{build/init/au-init} and is later installed into
the root filesystem as \file{/sbin/init}, matching the \code{init=/sbin/init}
parameter in \file{grub.cfg}.

\begin{infobox}[Verifying the static link]
The init agent's workflow recommends confirming the link mode after every build
with \code{file build/init/au-init}, which should report ``statically linked''.
If it ever reports ``dynamically linked,'' the binary risks failing to start at
boot --- a regression worth catching in CI.
\end{infobox}

\section{Design Rationale and Limitations}

\file{au-init} is a focused, honest implementation of the core PID~1 duties:
mounting the virtual filesystems, bringing up loopback, reading the hostname,
parsing a simple service table, supervising children with crash classification
and exponential-backoff restarts, reaping orphans, and tearing the system down
cleanly on signal. Its single-file, dependency-free, statically linked design is
exactly what a robust PID~1 should be.

It is equally important to be clear about what it is \emph{not} yet:

\begin{warnbox}[Current limitations]
\begin{itemize}
  \item \textbf{No service dependencies or ordering.} Services in
        \file{services.conf} are started in file order with no notion of ``start
        B after A is ready.'' There is no readiness or health probing beyond
        ``is the process still alive.''
  \item \textbf{Empty default service set.} The shipped \file{services.conf} is
        empty, so out of the box the only supervised process is the login shell.
        Background services must be added manually.
  \item \textbf{No socket activation, cgroups, or namespaces.} Unlike
        systemd-class inits, \file{au-init} does not place services in cgroups,
        apply resource limits, or perform lazy socket activation.
  \item \textbf{No runtime control channel.} Services cannot be started,
        stopped, or queried at runtime; the only external controls are the
        reboot/power-off signals to PID~1.
  \item \textbf{No \file{/etc/fstab} processing.} Init mounts a fixed,
        hard-coded set of pseudo-filesystems; it does not read \file{/etc/fstab}
        to mount additional real filesystems.
  \item \textbf{Fixed limits.} At most \code{MAX\_SERVICES} (32) services and 63
        arguments per command are supported, and argument parsing is naive
        whitespace splitting with no quoting support.
\end{itemize}
\end{warnbox}

These are reasonable boundaries for a teaching-oriented, minimalist
distribution. The architecture --- a flat service table, a flag-setting signal
handler, and a polling supervision loop --- is simple enough to read in one
sitting yet complete enough to keep a real system alive, restart crashing
services sanely, and shut down without corrupting the disk. That balance of
clarity and correctness is the defining virtue of AULinux's init.
